% ============================================================
% CHAPITRE 2 : ARCHITECTURE LOGICIELLE ET FRAMEWORK DE DÉVELOPPEMENT
% ============================================================
\chapter{Architecture logicielle et framework de développement}

\section{Introduction}

Avant de plonger dans la réalisation des sprints, il est indispensable de présenter les fondations techniques sur lesquelles repose l'ensemble de la solution. Ce chapitre détaille l'environnement matériel et logiciel mobilisé, l'architecture de déploiement retenue ainsi que l'organisation interne des microservices et de leurs interconnexions.

% --------------------------------------------------
\section{Environnement de développement}

\subsection{Environnement matériel}

\begin{table}[H]
    \centering
    \caption{Poste de travail utilisé}
    \label{tab:env_materiel}
    \begin{tabular}{|l|l|}
        \hline
        \textbf{Composant} & \textbf{Spécification} \\
        \hline
        Machine & PC portable Lenovo \\
        \hline
        Processeur & Intel Core i7 13650HX \\
        \hline
        Mémoire vive & 16 Go RAM \\
        \hline
        Stockage & SSD 512 Go \\
        \hline
        Système d'exploitation & Windows 11  \\
        \hline
    \end{tabular}
\end{table}

\subsection{Environnement logiciel}

Le Tableau~\ref{tab:env_logiciel} recense les outils mobilisés.

\begin{table}[H]
    \centering
    \caption{Outils et technologies du projet}
    \label{tab:env_logiciel}
    \renewcommand{\arraystretch}{1.3}
    \begin{tabular}{|p{3.5cm}|p{3cm}|p{6cm}|}
        \hline
        \textbf{Outil / Technologie} & \textbf{Version} & \textbf{Rôle dans le projet} \\
        \hline
        Visual Studio Code & Dernière & Environnement de rédaction du code \\
        \hline
        Node.js & 20+ & Moteur d'exécution des microservices \\
        \hline
        TypeScript & 5.x & Couche de typage pour fiabiliser les échanges \\
        \hline
        React & 19 & Construction des écrans (client et admin) \\
        \hline
        Vite & 7.x & Compilation ultra-rapide du Frontend \\
        \hline
        Tailwind CSS & 4.x & Habillage graphique via classes utilitaires \\
        \hline
        Express.js & 4.x & Routage HTTP interne des API \\
        \hline
        Prisma ORM & 5.x & Pont de communication avec PostgreSQL \\
        \hline
        PostgreSQL & 16 & Stockage des menus, utilisateurs et commandes \\
        \hline
        Python & 3.14 & Scripts d'intelligence artificielle \\
        \hline
        FastAPI & Dernière & Interface rapide pour le chatbot \\
        \hline
        LangChain & 0.3.x & Orchestrateur des requêtes RAG \\
        \hline
        Google Gemini & 2.5 Flash & Modèle linguistique générant les réponses \\
        \hline
        Figma & En ligne & Maquettage graphique \\
        \hline
        Git & Dernière & Versionnement du code source \\
        \hline
        Postman & Dernière & Simulation et validation des requêtes API \\
        \hline
        Cloudinary & Cloud & Hébergement externe des photos de plats \\
        \hline
        FAISS (faiss-cpu) & 1.7+ & Index vectoriel pour le pipeline RAG \\
        \hline
    \end{tabular}
\end{table}

% --------------------------------------------------
\section{Architecture de déploiement}

Le diagramme de déploiement (figure~\ref{fig:deploiement}) détaille la distribution physique des composants logiciels de la plateforme Assiette Gala, en illustrant les nœuds d'exécution (poste client, serveur d'application, serveur de bases de données) et les services externes mobilisés (Cloudinary, Google Gemini, Flouci, SMTP).

\begin{figure}[H]
    \centering
    \includegraphics[width=0.95\textwidth]{deploiement.png}
    \caption{Diagramme de déploiement de la plateforme Assiette Gala}
    \label{fig:deploiement}
\end{figure}

La figure~\ref{fig:architecture_globale} offre une vue d'ensemble de l'architecture logicielle de la plateforme, mettant en évidence les flux de communication entre le frontend, la passerelle API, les microservices métier, le service IA et les bases de données.

\begin{figure}[H]
    \centering
    \includegraphics[width=0.95\textwidth]{architecture_globale.png}
    \caption{Architecture globale de la plateforme Assiette Gala}
    \label{fig:architecture_globale}
\end{figure}

% --------------------------------------------------
\section{Architecture technique détaillée}

\subsection{Découpage en microservices}

La solution Assiette Gala se compose de sept briques logicielles autonomes, chacune spécialisée dans un domaine métier :

\begin{table}[H]
    \centering
    \caption{Microservices de la plateforme Assiette Gala}
    \label{tab:microservices}
    \renewcommand{\arraystretch}{1.3}
    \begin{tabular}{|p{3.5cm}|c|p{6cm}|}
        \hline
        \textbf{Service} & \textbf{Port} & \textbf{Responsabilité} \\
        \hline
        API Gateway & 3000 & Point d'entrée unique, routage, authentification, WebSocket \\
        \hline
        Auth Service & 3001 & Inscription, connexion, JWT, Google OAuth, profils \\
        \hline
        Catalog Service & 3002 & Plats, catégories, allergènes, stock, fournisseurs \\
        \hline
        Order Service & 3003 & Commandes, paiements, factures PDF \\
        \hline
        Event Service & 3004 & Événements, devis, modèles de réception \\
        \hline
        AI Service & 8000 & Chatbot RAG, recommandations, prévisions, cuisine \\
        \hline
        Frontend (Vite) & 5174 & Interface React (client et administration) \\
        \hline
    \end{tabular}
\end{table}

\subsection{Passerelle d'entrée (API Gateway)}

Pour éviter que le navigateur du client n'ait à connaître l'adresse de sept serveurs distincts, une \textbf{passerelle unique} (port 3000) joue le rôle de standard téléphonique : elle réceptionne chaque appel HTTP, vérifie le badge JWT du demandeur, puis achemine la requête vers la bonne cuisine logicielle.

\subsection{Package partagé (\texttt{@traiteurpro/shared})}

Afin d'éliminer la duplication de code entre les microservices, un package npm interne centralise les utilitaires communs :

\begin{itemize}
    \item \textbf{Middleware d'authentification :} vérification JWT, extraction du rôle utilisateur.
    \item \textbf{Gestionnaire d'erreurs :} format de réponse uniforme pour toutes les API.
    \item \textbf{Utilitaires de facturation :} génération des numéros de facture (format \texttt{INV-YYYYMM-<hex>}) et de commande (format \texttt{ORD-YYMMDD-<hex>}).
    \item \textbf{Types partagés :} interfaces TypeScript pour les entités transversales (utilisateur, commande, plat).
    \item \textbf{Validations :} schémas Zod pour la validation des entrées.
\end{itemize}

\subsection{Communication inter-services}

Les microservices communiquent entre eux via des appels HTTP internes sécurisés par un \textbf{jeton de service interne} (\texttt{INTERNAL\_SERVICE\_TOKEN}). Ce jeton, partagé exclusivement entre les services backend, empêche tout accès non autorisé aux endpoints internes.

Pour les notifications en temps réel, la passerelle API expose un serveur \textbf{WebSocket} (Socket.IO) qui diffuse les événements (nouvelle commande, mise à jour de statut) vers les clients connectés.

% --------------------------------------------------
\section*{Conclusion}
\addcontentsline{toc}{section}{Conclusion}

\begin{samepage}
Ce chapitre a dressé le panorama complet de l'outillage et de l'architecture technique de la plateforme Assiette Gala. Le découpage en microservices, la centralisation des utilitaires via le package partagé et la passerelle d'entrée unique garantissent une base solide pour les six sprints à venir. Le chapitre suivant inaugure le premier cycle de développement : la mise en place du système d'authentification.
\end{samepage}

